\chapter{§1.8 比与比例}
\label{sec:1.8}


\prerequisite{§1.5 分数，§1.7 百分数}

\gradelevel{6 年级}


\section*{比的意义与化简}


\begin{explore}


学校篮球队有 $12$ 名男生和 $8$ 名女生。老师说："男生和女生的人数比是 $12$ 比 $8$ 。"这个"比"是什么意思？它和除法、分数有什么关系？

在日常生活中，我们经常用"比"来描述两个量之间的关系。比如：

\begin{itemize}
  \item 调制蜂蜜水，蜂蜜和水的比例是 $1 : 5$
  \item 照片的长宽比是 $4 : 3$
  \item 混凝土中水泥、沙子、石子的比是 $1 : 2 : 3$
\end{itemize}


\end{explore}

\begin{understand}


\textbf{比}表示两个数量之间的倍数关系。 $a$ 和 $b$ 的比记作 $a : b$ ，读作" $a$ 比 $b$ "。

\[
a : b = \frac{a}{b}
\]


其中 $a$ 叫做比的\textbf{前项}， $b$ 叫做比的\textbf{后项}， $\frac{a}{b}$ 叫做\textbf{比值}。后项不能为 $0$ 。

比、除法、分数的关系：

\begin{center}
\begin{tabular}{llll}
\toprule
 & 比 & 除法 & 分数 \\
\midrule
形式 & $a : b$ & $a \div b$ & $\frac{a}{b}$ \\
各部分名称 & 前项 : 后项 & 被除数 $\div$ 除数 & 分子 / 分母 \\
结果 & 比值 & 商 & 分数值 \\
\bottomrule
\end{tabular}
\end{center}


\textbf{比的化简}：和分数约分类似，比的前项和后项同时除以它们的最大公因数。

\[
12 : 8 = 3 : 2
\]


化简的目标是使前项和后项都是整数，且互质（最大公因数为 $1$ ）。

如果比的前项或后项是分数或小数，先化成整数再化简：

\[
0.6 : 0.9 = 6 : 9 = 2 : 3
\]


\[
\frac{1}{3} : \frac{1}{2} = \frac{1}{3} \times 6 : \frac{1}{2} \times 6 = 2 : 3
\]


（前后项同乘以分母的最小公倍数。）

\end{understand}

\begin{workedexamples}


\textbf{例 1.} 化简 $24 : 36$ 。

\textbf{解：} $24$ 和 $36$ 的最大公因数是 $12$ 。

\[
24 : 36 = (24 \div 12) : (36 \div 12) = 2 : 3
\]


\textbf{例 2.} 化简 $1.5 : 2.5$ 。

\textbf{解：} 先化为整数比（同乘 $10$ ）：

\[
1.5 : 2.5 = 15 : 25 = 3 : 5
\]


\textbf{例 3.} 化简 $\frac{2}{3} : \frac{5}{6}$ 。

\textbf{解：} 前后项同乘 $6$ （分母的最小公倍数）：

\[
\frac{2}{3} : \frac{5}{6} = \frac{2}{3} \times 6 : \frac{5}{6} \times 6 = 4 : 5
\]


\textbf{例 4.} 小明身高 $1.5$ 米，爸爸身高 $1.8$ 米。小明和爸爸身高的比是多少？比值是多少？

\textbf{解：}
\begin{itemize}
  \item 身高比： $1.5 : 1.8 = 15 : 18 = 5 : 6$ 。
  \item 比值： $\frac{5}{6}$ （小明的身高是爸爸的 $\frac{5}{6}$ ）。
\end{itemize}


\end{workedexamples}

\begin{keytakeaway}


\begin{itemize}
  \item 比表示两个量的倍数关系， $a : b = \frac{a}{b}$ 。
  \item 比的前后项同时乘或除以同一个非零数，比值不变（类似分数基本性质）。
  \item 化简的目标：前后项为互质的整数。
  \item 比、除法、分数是同一种关系的三种表达方式。
\end{itemize}


\end{keytakeaway}

\begin{practice}


\begin{enumerate}
  \item 化简 $45 : 30$ 。
  \item 化简 $0.4 : 1.2$ 。
  \item 化简 $\frac{3}{4} : \frac{5}{8}$ 。
  \item 一杯果汁用 $40$ 毫升浓缩液加 $160$ 毫升水配成。浓缩液和水的比是多少？浓缩液占果汁总量的几分之几？
  \item 求 $7 : 4$ 的比值。
\end{enumerate}


\end{practice}

\section*{比例的意义与基本性质}


\begin{explore}


一幅地图上， $2$ 厘米代表实际的 $100$ 千米， $4$ 厘米代表实际的 $200$ 千米。这两组对应关系可以写成：

\[
2 : 100 = 4 : 200
\]


这个等式就是一个\textbf{比例}。比例有什么性质？怎么用它来解题？

\end{explore}

\begin{understand}


\textbf{比例}：表示两个比相等的式子。

\[
a : b = c : d \quad \text{即} \quad \frac{a}{b} = \frac{c}{d}
\]


其中 $a$ 和 $d$ 叫做\textbf{外项}， $b$ 和 $c$ 叫做\textbf{内项}。

\textbf{比例的基本性质（交叉相乘）}：在比例 $a : b = c : d$ 中，

\[
a \times d = b \times c
\]


即两个外项的积等于两个内项的积。

这条性质可以用来\textbf{解比例}——已知比例中的三个量，求第四个量。

\end{understand}

\begin{workedexamples}


\textbf{例 1.} 判断 $6 : 9$ 和 $10 : 15$ 能否组成比例。

\textbf{解：} 化简： $6 : 9 = 2 : 3$ ， $10 : 15 = 2 : 3$ 。两个比相等，所以能组成比例： $6 : 9 = 10 : 15$ 。

验证：外项之积 $= 6 \times 15 = 90$ ，内项之积 $= 9 \times 10 = 90$ 。相等，成立。

\textbf{例 2.} 解比例： $3 : 8 = x : 24$ 。

\textbf{解：} 由比例的基本性质：

\[
3 \times 24 = 8 \times x
\]


\[
72 = 8x
\]


\[
x = 9
\]


\textbf{例 3.} 解比例： $\frac{5}{x} = \frac{15}{12}$ 。

\textbf{解：} 交叉相乘：

\[
5 \times 12 = x \times 15
\]


\[
60 = 15x
\]


\[
x = 4
\]


\textbf{例 4.} 在一幅地图上， $3$ 厘米代表实际距离 $150$ 千米。如果两地在地图上相距 $7$ 厘米，实际距离是多少千米？

\textbf{解：} 设实际距离为 $x$ 千米。

\[
3 : 150 = 7 : x
\]


\[
3x = 150 \times 7 = 1050
\]


\[
x = 350
\]


两地实际距离为 $350$ 千米。

\end{workedexamples}

\begin{keytakeaway}


\begin{itemize}
  \item 比例是两个比相等的等式。
  \item 比例的基本性质：外项之积 $=$ 内项之积（交叉相乘）。
  \item 解比例的核心方法：交叉相乘，转化为乘法等式。
  \item 比例广泛应用于地图、缩放、配比等问题。
\end{itemize}


\end{keytakeaway}

\begin{practice}


\begin{enumerate}
  \item 判断 $4 : 7$ 和 $12 : 21$ 能否组成比例。
  \item 解比例： $x : 6 = 5 : 3$ 。
  \item 解比例： $\frac{2}{5} = \frac{x}{20}$ 。
  \item 解比例： $0.4 : 1.5 = x : 6$ 。
  \item 某地图上 $5$ 厘米代表实际距离 $200$ 千米。两城市实际相距 $560$ 千米，在地图上是多少厘米？
\end{enumerate}


\end{practice}

\section*{正比例与反比例}


\begin{explore}


情境一：每本笔记本 $3$ 元。买 $2$ 本要 $6$ 元，买 $5$ 本要 $15$ 元。本数越多，总价越高。总价和本数之间有什么固定的关系？

情境二：从家到学校 $6$ 千米。走路速度 $6$ 千米/时需 $1$ 小时；骑车速度 $12$ 千米/时需 $0.5$ 小时。速度越快，时间越短。速度和时间之间又有什么固定的关系？

\end{explore}

\begin{understand}


\textbf{正比例关系}：两个量 $x$ 和 $y$ ，如果它们的比值（商）始终不变，就说 $y$ 和 $x$ 成\textbf{正比例}。

\[
\frac{y}{x} = k \quad (\text{常数},\; k \neq 0)
\]


即 $y = kx$ 。一个量增大，另一个量也增大；一个量减小，另一个量也减小。

上面的情境一： $\frac{\text{总价}}{\text{本数}} = 3$ （元/本），总价和本数成正比例。

\textbf{反比例关系}：两个量 $x$ 和 $y$ ，如果它们的积始终不变，就说 $y$ 和 $x$ 成\textbf{反比例}。

\[
x \times y = k \quad (\text{常数},\; k \neq 0)
\]


即 $y = \frac{k}{x}$ 。一个量增大，另一个量反而减小。

上面的情境二： $\text{速度} \times \text{时间} = 6$ （千米），速度和时间成反比例。

\textbf{判断方法}：
\begin{enumerate}
  \item 找到两个变化的量。
  \item 如果它们的商不变——正比例。
  \item 如果它们的积不变——反比例。
  \item 如果商和积都在变——既不成正比例也不成反比例。
\end{enumerate}


\end{understand}

\begin{workedexamples}


\textbf{例 1.} 下面各组量是否成正比例或反比例？

(a) 长方形的面积一定，长和宽。

\textbf{解：} 长 $\times$ 宽 $=$ 面积（常数），所以长和宽成\textbf{反比例}。

(b) 正方形的边长和周长。

\textbf{解：} $\frac{\text{周长}}{\text{边长}} = 4$ （常数），所以周长和边长成\textbf{正比例}。

(c) 人的年龄和身高。

\textbf{解：} 年龄增大身高不一定按固定比例增长（长大后身高不再变），既不是正比例也不是反比例。

\textbf{例 2.} 一辆汽车从甲地开往乙地，速度和时间如下表：

\begin{center}
\begin{tabular}{lllll}
\toprule
速度（千米/时） & $40$ & $60$ & $80$ & $120$ \\
\midrule
时间（小时） & $6$ & $4$ & $3$ & $2$ \\
\bottomrule
\end{tabular}
\end{center}


判断速度和时间的关系。

\textbf{解：} 计算每组的积： $40 \times 6 = 240$ ， $60 \times 4 = 240$ ， $80 \times 3 = 240$ ， $120 \times 2 = 240$ 。积始终为 $240$ ，所以速度和时间成反比例。（ $240$ 千米就是甲乙两地之间的距离。）

\textbf{例 3.} 已知 $y$ 和 $x$ 成正比例，且当 $x = 4$ 时 $y = 10$ 。求当 $x = 6$ 时 $y$ 的值。

\textbf{解：} 因为成正比例，所以 $\frac{y}{x}$ 为常数。

\[
\frac{y}{x} = \frac{10}{4} = \frac{5}{2}
\]


当 $x = 6$ 时： $y = \frac{5}{2} \times 6 = 15$ 。

\textbf{例 4.} 已知 $y$ 和 $x$ 成反比例，且当 $x = 3$ 时 $y = 8$ 。求当 $x = 6$ 时 $y$ 的值。

\textbf{解：} 因为成反比例，所以 $x \times y$ 为常数。

\[
x \times y = 3 \times 8 = 24
\]


当 $x = 6$ 时： $y = \frac{24}{6} = 4$ 。

\end{workedexamples}

\begin{keytakeaway}


\begin{itemize}
  \item 正比例： $\frac{y}{x} = k$ （商不变），两个量同增同减。
  \item 反比例： $x \times y = k$ （积不变），一增一减。
  \item 判断关键：看"商"不变还是"积"不变。
  \item 不是所有两个相关的量都成比例关系。
\end{itemize}


\end{keytakeaway}

\begin{practice}


\begin{enumerate}
  \item 判断下面各组量是正比例、反比例还是不成比例：
\end{enumerate}

   (a) 每天工作量一定，完成任务所需天数和总工作量。
   (b) 路程一定，速度和时间。
   (c) 圆的面积和半径。
\begin{enumerate}
  \item $y$ 和 $x$ 成正比例，当 $x = 5$ 时 $y = 20$ 。求 $x = 8$ 时 $y$ 的值。
  \item $y$ 和 $x$ 成反比例，当 $x = 4$ 时 $y = 15$ 。求 $x = 10$ 时 $y$ 的值。
  \item 一项工程， $6$ 个人做需要 $10$ 天完成。如果增加到 $15$ 个人，需要几天？
  \item 购买钢笔，单价 $8$ 元。买 $5$ 支要多少钱？买 $12$ 支呢？总价和支数成什么比例？
\end{enumerate}


\end{practice}

\section*{比例的应用}


\begin{explore}


学校秋游要带 $120$ 瓶矿泉水，按照三个年级的人数 $2 : 3 : 5$ 来分配。每个年级各分到多少瓶？

另一个场景：一幅建筑设计图的比例尺是 $1 : 200$ ，图上量得一面墙长 $5$ 厘米。这面墙的实际长度是多少？

\end{explore}

\begin{understand}


#### 按比分配

把一个量按照给定的比来分配。

方法：先求出总份数，再用分数乘法分别计算。

以 $120$ 瓶水按 $2 : 3 : 5$ 分配为例：
\begin{itemize}
  \item 总份数 $= 2 + 3 + 5 = 10$
  \item 第一组： $120 \times \frac{2}{10} = 24$ （瓶）
  \item 第二组： $120 \times \frac{3}{10} = 36$ （瓶）
  \item 第三组： $120 \times \frac{5}{10} = 60$ （瓶）
\end{itemize}


验算： $24 + 36 + 60 = 120$ ，正确。

#### 比例尺

\textbf{比例尺}表示图上距离和实际距离的比。

\[
\text{比例尺} = \frac{\text{图上距离}}{\text{实际距离}}
\]


比例尺 $1 : 200$ 意味着图上 $1$ 厘米代表实际 $200$ 厘米（ $2$ 米）。

\[
\text{实际距离} = \frac{\text{图上距离}}{\text{比例尺}} = \text{图上距离} \times \frac{1}{\text{比例尺的比值}}
\]


\end{understand}

\begin{workedexamples}


\textbf{例 1.} 学校把 $540$ 本新书按 $4 : 3 : 2$ 分配给三个年级。每个年级各分到多少本？

\textbf{解：} 总份数 $= 4 + 3 + 2 = 9$ 。
\begin{itemize}
  \item 第一年级： $540 \times \frac{4}{9} = 240$ （本）。
  \item 第二年级： $540 \times \frac{3}{9} = 180$ （本）。
  \item 第三年级： $540 \times \frac{2}{9} = 120$ （本）。
\end{itemize}


验算： $240 + 180 + 120 = 540$ ，正确。

\textbf{例 2.} 一种混凝土由水泥、沙子、石子按 $1 : 2 : 4$ 混合而成。要配制 $2100$ 千克混凝土，各需多少千克？

\textbf{解：} 总份数 $= 1 + 2 + 4 = 7$ 。
\begin{itemize}
  \item 水泥： $2100 \times \frac{1}{7} = 300$ （千克）。
  \item 沙子： $2100 \times \frac{2}{7} = 600$ （千克）。
  \item 石子： $2100 \times \frac{4}{7} = 1200$ （千克）。
\end{itemize}


\textbf{例 3.} 一幅地图的比例尺是 $1 : 5000000$ （即 $1 : 500$ 万）。图上量得北京到天津的距离为 $2.4$ 厘米。北京到天津的实际距离约为多少千米？

\textbf{解：}

\[
\text{实际距离} = 2.4 \times 5000000 = 12000000 \text{（厘米）} = 120 \text{（千米）}
\]


北京到天津的实际距离约为 $120$ 千米。

\textbf{例 4.} 一栋教学楼实际高度为 $15$ 米。在一张比例尺为 $1 : 500$ 的设计图上，这栋楼画多高？

\textbf{解：} 先统一单位。 $15$ 米 $= 1500$ 厘米。

\[
\text{图上距离} = 1500 \times \frac{1}{500} = 3 \text{（厘米）}
\]


图上画 $3$ 厘米。

\textbf{例 5.} 甲乙两人合开一个小店，甲投资 $30000$ 元，乙投资 $20000$ 元。年底盈利 $15000$ 元，按投资比例分配。各得多少？

\textbf{解：} 投资比 $= 30000 : 20000 = 3 : 2$ ，总份数 $= 5$ 。
\begin{itemize}
  \item 甲： $15000 \times \frac{3}{5} = 9000$ （元）。
  \item 乙： $15000 \times \frac{2}{5} = 6000$ （元）。
\end{itemize}


\end{workedexamples}

\begin{keytakeaway}


\begin{itemize}
  \item 按比分配：先求总份数，再按分数乘法分配。
  \item 比例尺 $=$ 图上距离 $:$ 实际距离，注意统一单位。
  \item 实际距离 $=$ 图上距离 $\div$ 比例尺。
  \item 分配后一定要验算：各部分之和是否等于总量。
\end{itemize}


\end{keytakeaway}

\begin{practice}


\begin{enumerate}
  \item 把 $180$ 按 $2 : 3 : 4$ 分配。
  \item 三角形三个内角的度数比是 $1 : 2 : 3$ ，求三个角各是多少度？
  \item 一幅地图的比例尺是 $1 : 4000000$ 。图上 $3.5$ 厘米代表实际多少千米？
  \item 实际距离 $80$ 千米，在比例尺 $1 : 2000000$ 的地图上应画多少厘米？
  \item 甲乙丙三人合伙投资，甲出 $5$ 万，乙出 $3$ 万，丙出 $2$ 万。年底利润 $8$ 万元按投资比例分配，各得多少？
  \item 稀释一种农药，药液和水的比是 $1 : 150$ 。现有 $3$ 升药液，需要加多少升水？
  \item 制作一个学校操场的平面图。操场实际长 $200$ 米、宽 $100$ 米。如果用比例尺 $1 : 2000$ ，图上长和宽各画多少厘米？
  \item 甲仓库有 $300$ 吨货物，乙仓库有 $200$ 吨。从甲仓库调多少吨到乙仓库，才能使两个仓库的货物量之比为 $3 : 7$ ？
\end{enumerate}


\end{practice}



\begin{answer}


\textbf{练一练 比的意义与化简}

\begin{enumerate}
  \item $45 : 30 = 3 : 2$ 。
  \item $0.4 : 1.2 = 4 : 12 = 1 : 3$ 。
  \item $\frac{3}{4} : \frac{5}{8} = \frac{3}{4} \times 8 : \frac{5}{8} \times 8 = 6 : 5$ 。
  \item 浓缩液和水的比 $= 40 : 160 = 1 : 4$ 。果汁总量 $= 40 + 160 = 200$ （毫升），浓缩液占 $\frac{40}{200} = \frac{1}{5}$ 。
  \item $7 : 4$ 的比值 $= \frac{7}{4} = 1.75$ 。
\end{enumerate}


\textbf{练一练 比例的意义与基本性质}

\begin{enumerate}
  \item $4 : 7 = \frac{4}{7}$ ， $12 : 21 = \frac{12}{21} = \frac{4}{7}$ 。比值相等，能组成比例。
  \item $x \times 3 = 6 \times 5$ ， $3x = 30$ ， $x = 10$ 。
  \item $2 \times 20 = 5 \times x$ ， $40 = 5x$ ， $x = 8$ 。
  \item $0.4 \times 6 = 1.5 \times x$ ， $2.4 = 1.5x$ ， $x = 1.6$ 。
  \item 设图上距离为 $x$ 厘米。 $\frac{x}{5} = \frac{560}{200}$ ， $200x = 5 \times 560 = 2800$ ， $x = 14$ （厘米）。
\end{enumerate}


\textbf{练一练 正比例与反比例}

\begin{enumerate}
  \item (a) $\frac{\text{总工作量}}{\text{天数}} = \text{每天工作量}$ （常数），总工作量和天数成\textbf{正比例}。(b) 速度 $\times$ 时间 $=$ 路程（常数），成\textbf{反比例}。(c) 面积 $= \pi r^2$ ， $\frac{\text{面积}}{r}$ 不是常数， $\text{面积} \times r$ 也不是常数，\textbf{不成比例}。
  \item $\frac{y}{x} = \frac{20}{5} = 4$ 。当 $x = 8$ 时， $y = 4 \times 8 = 32$ 。
  \item $x \times y = 4 \times 15 = 60$ 。当 $x = 10$ 时， $y = \frac{60}{10} = 6$ 。
  \item 工作总量 $= 6 \times 10 = 60$ （人天）。 $15$ 人做需要 $60 \div 15 = 4$ （天）。
  \item $5$ 支： $8 \times 5 = 40$ （元）。 $12$ 支： $8 \times 12 = 96$ （元）。 $\frac{\text{总价}}{\text{支数}} = 8$ （常数），总价和支数成正比例。
\end{enumerate}


\textbf{练一练 比例的应用}

\begin{enumerate}
  \item 总份数 $= 2 + 3 + 4 = 9$ 。分别为 $180 \times \frac{2}{9} = 40$ ， $180 \times \frac{3}{9} = 60$ ， $180 \times \frac{4}{9} = 80$ 。验算： $40 + 60 + 80 = 180$ 。
  \item 三角形内角和 $= 180°$ ，总份数 $= 1 + 2 + 3 = 6$ 。三个角分别为 $30°$ 、 $60°$ 、 $90°$ 。
  \item $3.5 \times 4000000 = 14000000$ （厘米） $= 140$ （千米）。
  \item $80$ 千米 $= 8000000$ 厘米。 $8000000 \div 2000000 = 4$ （厘米）。
  \item 投资比 $= 5 : 3 : 2$ ，总份数 $= 10$ 。甲： $8 \times \frac{5}{10} = 4$ （万元），乙： $8 \times \frac{3}{10} = 2.4$ （万元），丙： $8 \times \frac{2}{10} = 1.6$ （万元）。
  \item $3 \times 150 = 450$ （升），需要加 $450$ 升水。
  \item $200$ 米 $= 20000$ 厘米， $20000 \div 2000 = 10$ （厘米）。 $100$ 米 $= 10000$ 厘米， $10000 \div 2000 = 5$ （厘米）。图上长 $10$ 厘米，宽 $5$ 厘米。
  \item 总货物 $= 300 + 200 = 500$ （吨）。按 $3 : 7$ 分配：甲 $= 500 \times \frac{3}{10} = 150$ （吨），乙 $= 500 \times \frac{7}{10} = 350$ （吨）。甲需调出 $300 - 150 = 150$ （吨）到乙。
\end{enumerate}


\end{answer}
